% ЗГЕНЕРОВАНО scripts/build_format_v2_chapter.py — не редагувати вручну.
% Джерело: editorial/cards/3.*.json, format-v2/rozdil-3.json, paraleli.json.

\naukchapter{3}{14}{Глава 3. Іспит пізнання}
\naukintro{Третя глава — найнеспокійніша з перших трьох. Вона не описує, як здобути знання, а показує, що знання саме по собі не вирішує нічого: у ній пʼять разів звучить aho там, де мудрець ловить себе на тому, що лишилося. Звідси — три останні вірші, де спокій описано не як здобуток, а як те, після чого питання про здобуток відпадає.}
\vsec{3.1}
\vdeva{अविनाशिनम् आत्मानम् एकं विज्ञाय तत्त्वतः । तवात्मज्ञस्य धीरस्य कथम् अर्थार्जने रतिः ॥}
\vtranslit{авінашінам атманам екам віджняя таттватах | таватмаджняся дгірася катхам артхарджане ратіх ||}
\vperek{Пізнавши істинно Атман — незнищенний і єдиний — як у тебе, атмазнавця й мудрого, може бути втіха в набутті користі?}
\vcom{Питання katham («як?») риторичне й будується на несумісності двох занять: пізнання єдиного й накопичення множинного. Слово artha тут не «сенс», а саме «користь, вигода» — це той самий термін, що в 10.1 серед того, що відкидають. Зверни увагу на порядок означень: спершу те, що пізнано (незнищенне, єдине), потім той, хто пізнав (атмазнавець, мудрий). Втіха в набутті не заборонена — поставлено питання, чи вона взагалі можлива за такого пізнання.}
\vpar{Українська паралель (наша інтерпретація)}{Григорій Сковорода, «Розмова пʼяти подорожніх про істинне щастя в житті».}{Наша паралель: у Сковороди щастя теж не лежить у набутку, і саме тому його не можна ні відняти, ні докупити. Наше порівняння — типологічне.}
\vfig{ilustratsii/3-1-balance.pdf}{Дві шальки: набуток і пізнаний Атман}{Схема наша, не історична ілюстрація.}

\vsec{3.2}
\vdeva{आत्माज्ञानाद् अहो प्रीतिर् विषयभ्रमगोचरे । शुक्तेर् अज्ञानतो लोभो यथा रजतविभ्रमे ॥}
\vtranslit{атмаджнянад ахо прітір вішаябграмаґочаре | шуктер аджнянато лобго ятха раджатавібграме ||}
\vperek{О диво! Через незнання Атмана з'являється прихильність у полі омани чуттєвих предметів — як через незнання перламутру з'являється жадоба в омані срібла.}
\vcom{Вірш будується на точній паралелі з двох частин, і паралель тут не прикраса, а аргумент: prīti (прихильність) співвіднесена з lobha (жадобою), viṣayabhrama (омана чуттєвого) — з rajatavibhrama (омана срібла), ātmājñāna (незнання Атмана) — з śuktyajñāna (незнання перламутру). Причина в обох випадках одна — незнання, а не вада предмета: перламутр не винен, що його беруть за срібло. Це та сама паралель, що в 2.9, але там був перелік трьох прикладів, а тут — один, розгорнутий до кінця. Слово prīti м'якше за spṛhā (2.22) і за lobha: це радше «прихильність», ніж «спрага».}
\vpar{Українська паралель (ПОТРЕБУЄ ПЕРЕВІРКИ)}{українська барокова проповідь, не визначено.}{ПОТРЕБУЄ ПЕРЕВІРКИ: перламутр і срібло як пара омани близькі до барокової емблематики, але конкретного твору ми не звірили й тому не називаємо.}
\vfig{ilustratsii/3-2-illusion.pdf}{Перламутр і срібло: жадоба з незнання}{Схема наша, не історична ілюстрація.}

\vsec{3.3}
\vdeva{विश्वं स्फुरति यत्रेदं तरङ्गा इव सागरे । सो ऽहम् अस्मीति विज्ञाय किं दीन इव धावसि ॥}
\vtranslit{вішвам спхураті ятредам таранґа іва саґаре | со ’хам асміті віджняя кім діна іва дгавасі ||}
\vperek{Пізнавши: «я — те, у чому сяє весь цей світ, як хвилі в морі», — чого ти біжиш, наче вбогий?}
\vcom{Питання kim ... dhāvasi («чого біжиш») спрямоване не до дії взагалі, а до бігу як способу життя: дієслово dhāv означає саме фізичний біг, метушню. Порівняння з dīna («убогий») вказує на мотив, а не на стан: біжать ті, кому бракує. Формула so 'ham asmi («я є те») — та сама, що в 2.15, але там вона завершувала заперечення трійці, а тут відкриває питання. Образ хвиль у морі повторює 2.10 і 2.23–2.25: це наскрізний образ тексту, а не місцева метафора.}
\vpar{Українська паралель (ПОТРЕБУЄ ПЕРЕВІРКИ)}{українська філософська лірика, не визначено.}{ПОТРЕБУЄ ПЕРЕВІРКИ: докір тому, хто шукає те, що вже має, в українській ліриці трапляється, але назвати твір без звірки ми не можемо.}
\vfig{ilustratsii/3-3-motion.pdf}{Хвиля, що шукає моря, у якому вже стоїть}{Схема наша, не історична ілюстрація.}

\vsec{3.4}
\vdeva{श्रुत्वापि शुद्धचैतन्यम् आत्मानम् अतिसुन्दरम् । उपस्थे ऽत्यन्तसंसक्तो मालिन्यम् अधिगच्छति ॥}
\vtranslit{шрутвапі шуддгачайтаням атманам атісундарам | упастхе ’тянтасамсакто маліням адгіґаччаті ||}
\vperek{Навіть почувши про Атман — чисту свідомість, надзвичайно прекрасну, — той, хто надмірно прив'язаний до того, що під рукою, набуває бруду.}
\vcom{Слово upasthe навмисно не уточнене: воно означає і «на лоні», і «поблизу, під рукою», і традиційно вживається з плотським відтінком. Точний зміст залежить від контексту, якого вірш не дає, — і саме тому переклад зберігає двозначність, а не обирає одне значення. Структура вірша — контраст між почутим і робленим: śrutvā api («навіть почувши») підкреслює, що знання про Атман саме собою нічого не змінює. mālinya («бруд») — не моральна оцінка, а стан забрудненості; те саме слово вживане про одяг і про посуд. Тобто йдеться не про гріх, а про те, що річ втрачає чистоту від дотику.}
\vpar{Українська паралель (наша інтерпретація)}{Іван Франко, «Зівʼяле листя» (1896).}{Наша паралель: у Франка знання теж не рятує від привʼязаності — воно лише робить її виднішою. Наше порівняння — за способом бачити людину, а не за вченням.}
\vfig{ilustratsii/3-4-balance.pdf}{Дві шальки: почуте про Атман і те, що під рукою}{Схема наша, не історична ілюстрація.}

\vsec{3.5}
\vdeva{सर्वभूतेषु चात्मानं सर्वभूतानि चात्मनि । मुनेर् जानत आश्चर्यं ममत्वम् अनुवर्तते ॥}
\vtranslit{сарвабгутешу чатманам сарвабгутані чатмані | мунер джаната ашчар'ям маматвам анувартате ||}
\vperek{О диво: у мудреця, який знає Атман у всіх істотах і всі істоти — в Атмані, досі тримається «моє».}
\vcom{Двозмінна конструкція (Атман у всіх, усі в Атмані) — не повтор, а два напрямки одного знання: назовні й усередину. Диво в тому, що навіть за такого знання лишається mamatvam — не «моє» як майно, а мійність як звичка. Дієслово anuvartate («іде слідом») важливе: воно не каже, що мійність панує, а що вона тягнеться позаду. Це найточніше місце глави про межі знання: знання не вичищає звичку автоматично. Порівняй 3.4, де йшлося про того, хто лише почув: там бруд набувають, тут мійність лишається — різні ступені, і текст не згладжує різниці.}
\vpar{Українська паралель (ПОТРЕБУЄ ПЕРЕВІРКИ)}{українська релігійна думка, не визначено.}{ПОТРЕБУЄ ПЕРЕВІРКИ: докір мудрецеві, в якого лишилося своє, близький до української аскетичної літератури, але конкретного місця ми не перевірили.}
\vfig{ilustratsii/3-5-knot.pdf}{Вузол, що лишається навіть у знанні}{Схема наша, не історична ілюстрація.}

\vsec{3.6}
\vdeva{आस्थितः परमाद्वैतं मोक्षार्थे ऽपि व्यवस्थितः । आश्चर्यं कामवशगो विकलः केलिशिक्षया ॥}
\vtranslit{астхітах парамадвайтам мокшартхе ’пі в'явастхітах | ашчар'ям камавашаґо вікалах келішікшая ||}
\vperek{Той, хто утвердився у вищій недвоїстості й стоїть навіть у справі звільнення, — о диво! — під владою бажання: безпорадний, мов у грі-вправлянні.}
\vcom{Вірш починає серію із чотирьох «див» (3.6–3.8, далі 3.10), і кожне будується за однією схемою: високий стан — і те, що в ньому несподівано лишається. Тут протиставлені два утвердження (āsthitaḥ, vyavasthitaḥ — майже синоніми, поставлені поруч навмисно) і один стан, що їх перекреслює. Слово vikala («безпорадний») сильніше за просте «слабкий»: воно означає позбавленість частини, скаліченість. keliśikṣayā («грою-вправлянням») знижує регістр: те, що мало бути шляхом, виявилося забавкою.}
\vpar{Українська паралель (наша інтерпретація)}{Василь Стус, «Веселий цвинтар» (1985).}{Наша паралель: у Стуса людина теж стоїть на своєму — і все одно бажання веде її далі, ніж вона хотіла. Наша асоціація, не твердження про вплив.}
\vfig{ilustratsii/3-6-knot.pdf}{Вузол бажання, затягнутий поверх недвоїстості}{Схема наша, не історична ілюстрація.}

\vsec{3.7}
\vdeva{उद्भूतं ज्ञानदुर्मित्रम् अवधार्यातिदुर्बलः । आश्चर्यं कामम् आकाङ्क्षेत् कालम् अन्तम् अनुश्रितः ॥}
\vtranslit{удбгутам джнянадурмітрам авадгар'ятідурбалах | ашчар'ям камам аканкшет калам антам анушрітах ||}
\vperek{О диво! Зрозумівши, що знання виявилося поганим другом, — надзвичайно слабкий — він усе одно прагнув би бажання, дійшовши до кінця часу.}
\vcom{Складене jñānadurmitram можна прочитати двома способами: «знання, що є поганим другом» і «поганий друг, [який полягає] у знанні». Обидва дають той самий сенс: знання, на яке покладалися, не врятувало. Найсильніше тут — час дієслова й обставина: ākāṅkṣet стоїть в optativus («прагнув би»), а kālam antam anuśritaḥ («дійшовши до кінця часу») означає останню межу, коли вже нічого не змінити. Тобто бажання не просто лишається — воно лишається там, де його вже неможливо задовольнити. Це не докір, а констатація абсурду, і слово āścaryam («диво») вжите саме в цьому сенсі.}
\vpar{Українська паралель (ПОТРЕБУЄ ПЕРЕВІРКИ)}{українська поезія XX ст., не визначено.}{ПОТРЕБУЄ ПЕРЕВІРКИ: мотив знання, що виявилося поганим другом, близький до української лірики про зраду власних переконань, але конкретного твору ми не називаємо без звірки.}
\vfig{ilustratsii/3-7-flame.pdf}{Полумʼя, що не зігріває: знання як поганий друг}{Схема наша, не історична ілюстрація.}

\vsec{3.8}
\vdeva{इहामुत्र विरक्तस्य नित्यानित्यविवेकिनः । आश्चर्यं मोक्षकामस्य मोक्षाद् एव विभीषिका ॥}
\vtranslit{іхамутра вірактася нітянітявівекінах | ашчар'ям мокшакамася мокшад ева вібгішіка ||}
\vperek{О диво: у того, хто відвернувся від цього й того світу, хто розрізняє вічне й невічне, хто бажає звільнення, — страх саме від звільнення.}
\vcom{Найгостріший вірш серії. Три ознаки поставлені поспіль і всі три «правильні»: відречення, розрізнення, бажання звільнення. І саме з третьої виростає страх, бо звільнення означає втрату того, хто звільняється. ablativus mokṣāt («від звільнення») з eva («саме») робить джерелом страху не світ і не смерть, а мету. vibhīṣikā — не «боязкість», а те, що лякає; слово вживане і про опудало. Тобто звільнення постає як те, чим лякають.}
\vpar{Українська паралель (наша інтерпретація)}{Леся Українка, «Кассандра» (1907).}{Наша паралель: у Лесі Українки знання теж не дає спокою, а забирає його, — і страх приходить саме від того, що має визволити. Наше порівняння — типологічне.}
\vfig{ilustratsii/3-8-knot.pdf}{Вузол, завʼязаний на самому звільненні}{Схема наша, не історична ілюстрація.}

\vsec{3.9}
\vdeva{धीरस् तु भोज्यमानो ऽपि पीड्यमानो ऽपि सर्वदा । आत्मानं केवलं पश्यन् न तुष्यति न कुप्यति ॥}
\vtranslit{дгірас ту бгоджямано ’пі підямано ’пі сарвада | атманам кевалам пашян на тушяті на куп'яті ||}
\vperek{Мудрий же — чи його голублять, чи його мучать — завжди, бачачи лише Атман, і не радіє, і не гнівається.}
\vcom{Два дієприкметники пасиву (bhojyamāna, pīḍyamāna) поставлені як рівноцінні умови: похвала й мука описані однією граматичною формою, тобто для мудрого вони вже не різні події. Це не байдужість: na tuṣyati na kupyati — два окремі заперечення, а не одне «не хвилюється». Слово kevalam («лише») робить ключову роботу: бачить не «Атман крізь усе», а лише Атман, і саме тому немає кому радіти. Форма дієприкметника paśyan (теперішній час) вказує на тривалість, а не на одноразове прозріння.}
\vpar{Українська паралель (наша інтерпретація)}{Тарас Шевченко, «Неофіти» (1857).}{Наша паралель: у Шевченка гідність теж не залежить від того, хвалять людину чи мучать. Наше порівняння — за будовою думки, не про джерело.}
\vfig{ilustratsii/3-9-balance.pdf}{Дві шальки: хвала й мука — важать однаково}{Схема наша, не історична ілюстрація.}

\vsec{3.10}
\vdeva{चेष्टमानं शरीरं स्वं पश्यत्य् अन्यशरीरवत् । संस्तवे चापि निन्दायां कथं क्षुभ्येत् महाशयः ॥}
\vtranslit{чештаманам шарірам свам пашятй аняшарірават | самставе чапі ніндаям катхам кшубгєт махашаях ||}
\vperek{Своє тіло, що рухається, він бачить як тіло чуже. То як у хвалі чи в ганьбі збентежився б великодухий?}
\vcom{Вірш закриває серію «див» і робить це тихо, без āścaryam: тут немає вигуку, є питання katham («як?»), на яке відповідь очевидна. ceṣṭamānam («те, що ворушиться») вибране навмисно: тіло описане не як річ, а як процес, і саме тому його легко побачити збоку. Пара saṃstava / nindā («хвала» / «ганьба») — та сама пара, що bhojyamāna / pīḍyamāna у 3.9, але там ішлося про дію на тіло, тут — про слово. mahāśayaḥ («великодухий») буквально означає «той, у кого велике ложе», і це не метафора розуму, а образ простору, у якому вміщається і хвала, і ганьба.}
\vpar{Українська паралель (наша інтерпретація)}{Олександр Довженко, «Щоденник» (1941–1956).}{Наша паралель: у Довженка тіло теж постає як щось, що людині належить менше, ніж їй здається. Наше порівняння — за способом бачити себе.}
\vfig{ilustratsii/3-10-witness.pdf}{Тіло як чуже в колі побаченого, бачитель у центрі}{Схема наша, не історична ілюстрація.}

\vsec{3.11}
\vdeva{मायामात्रम् इदं विश्वं पश्यन् विगतकौतुकः । अपि सन्निहिते मृत्यौ कथं त्रस्यति धीरधीः ॥}
\vtranslit{маяматрам ідам вішвам пашян віґатакаутуках | апі санніхіте мрітяу катхам трасяті дгірадгіх ||}
\vperek{Бачачи, що весь цей світ — тільки мая, і втративши до нього охоту, — навіть коли смерть уже поруч, як злякався б той, чий розум стійкий?}
\vcom{Слово kautuka тут важить більше, ніж здається: це не просто «цікавість», а охота до святкового, потяг до видовища. Тобто йдеться не про втрату знань, а про втрату інтересу до гри. locativus absolutus sannihite mṛtyau («коли смерть присутня») ставить умову максимально жорстко: не «якось колись», а тут і тепер. dhīradhīḥ («чий розум стійкий») — не «хоробрий»: стійкість розуму, а не безстрашність характеру. Питання katham знову риторичне, як у 3.3 і 3.10 — глава послідовно закриває питання питаннями.}
\vpar{Українська паралель (наша інтерпретація)}{Михайло Коцюбинський, «Тіні забутих предків» (1911).}{Наша паралель: у Коцюбинського світ теж постає як те, що показується й зникає, не питаючи людину. Наше порівняння — за образом, а не за вченням.}
\vfig{ilustratsii/3-11-illusion.pdf}{Серпанок маї: світ, що зникає під поглядом}{Схема наша, не історична ілюстрація.}

\vsec{3.12}
\vdeva{निःस्पृहं मानसं यस्य नैराश्ये ऽपि महात्मनः । तस्यात्मज्ञानतृप्तस्य तुलना केन जायते ॥}
\vtranslit{ніхспріхам манасам яся найрашє ’пі махатманах | тасятмаджнянатріптася тулана кена джаяте ||}
\vperek{У кого розум без прагнень, хто навіть у безнадії лишається великим духом, хто вдоволений знанням Атмана, — з ким його й порівнювати?}
\vcom{nairāśye 'pi («навіть у безнадії») — найтонше місце. Безнадія тут не недолік, а умова: відсутність сподівань і є тим, що робить розум вільним від прагнень. Тому й сказано api («навіть») — ніби всупереч очікуванню, що без надії людина меншає. tulanā («порівняння, міра») вжите як річ, що не виникає: не «він не має рівних», а «порівняння не народжується», тобто зникає сама операція зіставлення. Форма jāyate (теперішній, «народжується») показує це як постійний стан, а не як разовий висновок.}
\vpar{Українська паралель (ПОТРЕБУЄ ПЕРЕВІРКИ)}{українська аскетична проза, не визначено.}{ПОТРЕБУЄ ПЕРЕВІРКИ: опис розуму без прагнень, з яким немає з чим порівняти, близький до української аскетичної прози, але конкретного твору ми не перевірили.}
\vfig{ilustratsii/3-12-space.pdf}{Небо без хмар: розум без прагнень}{Схема наша, не історична ілюстрація.}

\vsec{3.13}
\vdeva{स्वभावाद् एव जानानो दृश्यम् एतन् न किञ्चन । इदं ग्राह्यम् इदं त्याज्यं स किं पश्यति धीरधीः ॥}
\vtranslit{свабгавад ева джанано дрішям етан на кіньчана | ідам ґрахям ідам тяджям са кім пашяті дгірадгіх ||}
\vperek{Знаючи самою природою, що все видиме — аніщо, — «це треба взяти, це треба відкинути» — то що ж бачить той, чий розум стійкий?}
\vcom{Дві пари герундивів (grāhyam / tyājyam) подані в лапках як чужа мова: це не позиція мудрого, а те, що він чує навколо або в собі. Заперечується не вибір між «взяти» й «відкинути», а сама наявність пари. svabhāvād eva («самою природою») знімає потребу в доказі: це не висновок, а спосіб бачення. Питання kim paśyati («що бачить») точніше за «чи бачить»: воно питає про предмет, а не про факт, бо коли видиме — аніщо, предмета не лишається.}
\vpar{Українська паралель (наша інтерпретація)}{Григорій Сковорода, «Байки харківські» (1774).}{Наша паралель: у Сковороди вибір між речами теж виявляється зайвим для того, хто вже бачить, що речі не важать. Наше порівняння — типологічне.}
\vfig{ilustratsii/3-13-neti.pdf}{Взяти й відкинути відпадають: усе видиме — аніщо}{Схема наша, не історична ілюстрація.}

\vsec{3.14}
\vdeva{अन्तस्त्यक्तकषायस्य निर्द्वन्द्वस्य निराशिषः । यदृच्छयागतो भोगो न दुःखाय न तुष्टये ॥}
\vtranslit{антастяктакашаяся нірдвандвася нірашішах | ядріччаяґато бгоґо на духкхая на туштає ||}
\vperek{У того, в кого всередині відкинуто каламуть, хто без двоїстості й без сподівань, усе, що випадає само собою, не стає ні стражданням, ні задоволенням.}
\vcom{Давальні duḥkhāya і tuṣṭaye — не «не приносить», а «не йде на»: досвід не зараховується ні туди, ні туди. Це тонше за байдужість: байдужість усе ще ставить оцінку «ніяк», а тут досвід просто не конвертується в жодну з двох категорій. kaṣāya («каламуть, осад») — слово з мови фарб і напоїв: це не гріх і не пристрасть у моральному сенсі, а те, що каламутить прозорість. antaḥ («усередині») стоїть першим членом навмисно: відкинуто не зовнішнє, а осад усередині. yadṛcchayāgataḥ («прийшов сам собою») — те саме yadṛcchā, що в 4.4 і далі; тут воно означає не випадковість, а відсутність замовлення.}
\vpar{Українська паралель (наша інтерпретація)}{Микола Зеров, «Камена» (1924).}{Наша паралель: у Зерова ясність теж приходить після того, як каламуть відстоялася, і тоді все, що випадає, приймають рівно. Наше порівняння — за образом, не за джерелом.}
\vfig{ilustratsii/3-14-space.pdf}{Прозоре небо: каламуть відстоялася}{Схема наша, не історична ілюстрація.}

